\section{Models and strategies}
\label{sec:strategies}

\subsection{Models}
\label{ssec:models}

We evaluate six production frontier models that span $\sim$20$\times$ in price and represent three API families (Anthropic, OpenAI, Google):

\begin{itemize}
\item \opus (Anthropic): \dollars{15}/\$75 per million input/output tokens.
\item \sonnet (Anthropic): \dollars{3}/\$15.
\item \haiku (Anthropic): \dollars{1}/\$5.
\item \gpt (OpenAI Responses API): \dollars{5}/\$20.
\item \geminipro (Google AI Studio API, paid tier, $\leq$200K context): \dollars{1.25}/\$10.
\item \geminiflash (Google AI Studio API, paid tier, $\leq$200K context): \dollars{0.30}/\$2.50.
\end{itemize}

All runs use temperature $= 1.0$ (default for both APIs).
Costs are computed using each model's published per-token rate as of April 2026; \cref{sec:results} reports both raw token counts and dollarised costs.

\subsection{Strategies}
\label{ssec:strategies}

\GeoGenBench ships two reference strategies as comparators, each implementing a common \texttt{SubstanceStrategy} interface.
The benchmark itself does not depend on either --- any text-emitting model can be benchmarked by routing its TikZ output through the verification pipeline.
The strategies sit on a spectrum from ``maximum LLM freedom'' (\rawcode) to ``schema-constrained intermediate representation'' (\structured); comparing them isolates the contribution of typed intermediate representations and pre-render self-validation.
\Cref{fig:strategies_compare} shows what the LLM emits in each strategy on the same prompt.

\paragraph{\rawcode (single-shot direct TikZ).}
The LLM is given the prompt plus a TikZ tutorial system prompt and produces TikZ via a \texttt{render\_diagram(tikz: str)} tool call.
It may call the tool multiple times to revise based on the rendered output, but no symbolic checks are run.
This is the zero-scaffolding baseline.

\paragraph{\structured (schema-constrained intermediate representation).}
The LLM is given a Pydantic schema for an intermediate \texttt{DiagramIR} representation (point declarations, segment declarations, intersections, circles, render ops, and \emph{check} declarations) and must emit IR JSON.
Code (not the LLM) compiles IR $\to$ SymPy objects $\to$ TikZ $\to$ SVG.
Validation against the IR's check list happens \emph{before} render, so the strategy can self-correct: if a declared check fails on the LLM's IR, the strategy revises (up to $N=3$ retries).
This is the strategy we recommend for production deployment.

\begin{figure}[t]
\centering
\small
% Common prompt header
\fbox{\begin{minipage}{0.97\linewidth}
\textit{Common prompt:} ``Draw an acute triangle $EFG$ with $\angle E = 60^{\circ}$, $\angle F = 70^{\circ}$. Construct the altitude from $G$ meeting $EF$ at $H$.''
\end{minipage}}\\[6pt]

% Two strategy columns. Each column = title + role line + code box of identical
% fixed height (so frames line up regardless of how many code lines there are).
\newlength{\stratcolwd}\setlength{\stratcolwd}{0.47\linewidth}
\newlength{\stratcodeht}\setlength{\stratcodeht}{52mm}

\begin{minipage}[t]{\stratcolwd}
\centering
{\textbf{\rawcode}}\\
{\scriptsize\itshape zero scaffolding}\\[2pt]
{\scriptsize LLM emits TikZ directly:}\\[2pt]
\fbox{\begin{minipage}[t][\stratcodeht][t]{\dimexpr\stratcolwd-2\fboxsep-2\fboxrule\relax}
\ttfamily\tiny
\textbackslash coordinate (E) at (0,0);\\
\textbackslash coordinate (F) at (5,0);\\
\textbackslash coordinate (G) at (3.07,5.31);\\
\textbackslash coordinate (H) at (3.07,0);\\
\textbackslash draw (E)--(F)--(G)--cycle;\\
\textbackslash draw[densely dashed]\\
\hphantom{xx}(G) -- (H);
\end{minipage}}
\end{minipage}\hfill
\begin{minipage}[t]{\stratcolwd}
\centering
{\textbf{\structured}}\\
{\scriptsize\itshape schema-constrained IR}\\[2pt]
{\scriptsize LLM emits typed JSON; \emph{we} compile to TikZ:}\\[2pt]
\fbox{\begin{minipage}[t][\stratcodeht][t]{\dimexpr\stratcolwd-2\fboxsep-2\fboxrule\relax}
\ttfamily\tiny
\{ "points": [\\
\hphantom{x}\{"name":"E","type":"free"\},\\
\hphantom{x}\{"name":"F","type":"free"\},\\
\hphantom{x}\{"name":"G","type":"by\_angles",\\
\hphantom{xx}"from":"E","to":"F",\\
\hphantom{xx}"ang\_at\_from":60,\\
\hphantom{xx}"ang\_at\_to":70\}\,],\\
\hphantom{x}"intersections":[\\
\hphantom{xx}\{"name":"H",\\
\hphantom{xxx}"lines":[\\
\hphantom{xxxx}["G","perp\_EF"],\\
\hphantom{xxxx}["E","F"]\,]\}\,],\\
\hphantom{x}"checks":[\{"type":\\
\hphantom{xxx}"right\_angle",\\
\hphantom{xxx}"args":["G","H","E"]\}\,]\}
\end{minipage}}
\end{minipage}

\vspace{4pt}
\centering\textit{\small$\Big\downarrow$\ \ Both converge to the same TikZ\,$\to$\,SVG\,$\to$\,SymPy\,$\to$\,predicate-verdict pipeline of \cref{fig:pipeline}.\ \ $\Big\downarrow$}

\caption{\textbf{The two reference generation strategies, side by side on the same prompt.}
\rawcode places minimal trust in scaffolding: the LLM emits coordinate-bearing TikZ directly.
\structured asks the LLM to declare \emph{what} it wants to draw via a typed schema (an intermediate representation, IR); \emph{our} compiler picks coordinates and the IR's own \texttt{checks} list lets the strategy self-correct before rendering.}
\label{fig:strategies_compare}
\end{figure}

\paragraph{Beyond TikZ: an alternate text-emit target.}
The methodology is renderer-agnostic --- the verifier does not care whether the model emits TikZ, raw SVG, Asymptote, or any other text-emittable vector format, provided we can compile it back to a coordinate set.
The codebase ships \texttt{strategies/raw\_svg.py} (and a self-revising variant) which routes the model directly to SVG output, bypassing TikZ and the LuaLaTeX container entirely.
Because raw SVG and raw TikZ measure essentially the same thing --- ``can the model emit a text-format vector diagram with correct geometry'' --- we do not include \texttt{raw\_svg} in the headline leaderboard; it is intended for deployments where a TeX-free runtime is required, and is documented in the released codebase rather than benchmarked here.

\subsection{Cost and latency methodology}
\label{ssec:cost_methodology}

For each (model, strategy, scenario, repeat) we record (a)~total input + output tokens across all API calls (including any retries), (b)~wall-clock duration from prompt submission to rendered SVG, and (c)~the token-multiplied USD cost at published rates.
The \emph{mean cost per scenario} reported in \cref{sec:results} is the arithmetic mean of per-scenario costs across the headline run, including failed scenarios (a model that frequently retries before failing pays for the retries).
Per-call latency excludes time spent in our verification pipeline.

\subsection{Headline-run configuration}
\label{ssec:pilot_config}

The headline run reported in \cref{sec:results} draws each of the $600$ templated scenarios with three repeats per (model, strategy) cell.
The \structured strategy permits up to three internal check-fail retries before recording final verdict.
Two cells are marginally partial --- \opus + \structured at $574/600$ unique scenarios with $\geq 3$ repeats and \gpt + \structured at $545/600$ --- because a fixed set of pathology scenarios trigger circular-dependency cycles in the IR resolver that the subprocess-isolated runner kills via OS-level timeout; we document the per-cell coverage caveat in \cref{ssec:headline} and the per-tier breakdown in \cref{ssec:tiers_results}.
